\chapter{Implémentation technique}
\label{chap:impl}

\section{Pile technologique}

Le tableau~\ref{tab:pile} présente la pile effectivement utilisée. Conformément au cadrage
MVP retenu, les versions et composants y sont indiqués tels qu'ils sont réellement mis en
œuvre ; les briques avancées non câblées sont reportées au chapitre~\ref{chap:perspectives}.

\begin{table}[H]
\centering
\caption{Pile technologique de BCaaS}
\label{tab:pile}
\renewcommand{\arraystretch}{1.25}\footnotesize
\begin{tabularx}{\linewidth}{L{4.2cm} X L{2.8cm}}
\toprule
\textbf{Composant} & \textbf{Technologie} & \textbf{Version} \\
\midrule
Framework backend     & Spring Boot (servlet, MVC)         & 3.5.x \\
Langage               & Java                               & 21 LTS \\
Base de données       & PostgreSQL                         & 16 \\
Migrations            & Liquibase                          & 4.x \\
Sécurité              & Spring Security + JWT (jjwt)       & 0.12.6 \\
Documentation API     & SpringDoc OpenAPI                  & 2.8.6 \\
Cache / pub-sub       & Redis                              & 7 \\
Bus d'événements      & Apache Kafka (infra)               & 7.6 (image) \\
Build                 & Maven                              & 3.9+ \\
Tests                 & JUnit 5 + Mockito + AssertJ        & 152 tests \\
Conteneurisation      & Docker + Docker Compose            & --- \\
Frontend              & Next.js / React / TypeScript       & 16 / 19 \\
\bottomrule
\end{tabularx}
\end{table}

\section{Structure du projet}

Le projet suit l'architecture hexagonale combinée au DDD. L'arborescence backend est
organisée par responsabilités techniques au sein du package \texttt{com.bizcore.bizcore\_backend}.

\begin{lstlisting}[language={},caption={Structure du projet backend},captionpos=b]
src/main/java/com/bizcore/bizcore_backend/
  +-- auth/         # Authentification (register, login, roles)
  +-- config/       # Configurations (Kafka, Redis, Security, cache)
  +-- controller/   # Controllers REST
  +-- domain/       # Entites JPA (15) + enum SupportedCurrency
  +-- dto/          # Objets de transfert
  +-- exception/    # GlobalExceptionHandler, ApiError
  +-- listener/     # Listener Kafka (consommateur)
  +-- repository/   # Repositories Spring Data (tenant-aware)
  +-- security/     # JwtService, JwtAuthFilter, TenantFilter, TenantContext
  +-- service/      # Services metier
src/main/resources/
  +-- application.properties
  +-- db/changelog/ # Migrations Liquibase
src/test/           # 16 classes, 152 tests (JUnit 5 + Mockito)
\end{lstlisting}

\section{Multi-tenant et isolation}

L'isolation repose sur le modèle « schéma partagé, colonne \texttt{tenant\_id} » et sur
trois pièces coordonnées.

\paragraph{1. \texttt{TenantContext} --- un \texttt{ThreadLocal}.}
Chaque requête HTTP dispose d'une copie isolée du tenant actif, stockée dans un
\texttt{ThreadLocal<UUID>}, ce qui évite toute fuite entre requêtes concurrentes partageant
le même pool de threads.

\paragraph{2. \texttt{TenantFilter} --- extraction depuis le JWT.}
Filtre \texttt{OncePerRequestFilter} d'ordre $2$ (après le filtre d'authentification), il
lit le claim \texttt{tenantId} du JWT et le place dans le contexte, puis le purge dans un
bloc \texttt{finally}.

\begin{lstlisting}[language=Java,caption={Extrait simplifié de TenantFilter},captionpos=b]
String tenantIdStr = jwtService.extractTenantId(jwt);   // claim "tenantId"
TenantContext.setTenantId(UUID.fromString(tenantIdStr));
try {
    filterChain.doFilter(request, response);
} finally {
    TenantContext.clear();   // pas de fuite inter-tenant
}
\end{lstlisting}

\paragraph{3. Injection du claim à la connexion.}
À l'authentification, \texttt{JwtService} embarque \texttt{tenantId} et \texttt{tenantName}
dans le token. Le tenant est ainsi dérivé de l'utilisateur (colonne \texttt{users.tenant\_id})
et relu à chaque requête sans appel supplémentaire à la base. L'isolation est appliquée
dans la couche service via des requêtes \emph{tenant-aware} de type
\texttt{findAllByTenantId(...)}.

\begin{garantie}[Garantie d'isolation et limite assumée]
Le \texttt{tenant\_id} circule dans le contexte d'exécution et conditionne les requêtes des
services principaux ; un générateur de clés de cache préfixe par tenant. \textbf{Limite
assumée :} l'isolation n'est pas une barrière globale (ni filtre Hibernate
\texttt{@Filter}, ni \emph{Row-Level Security} PostgreSQL) ; elle repose sur la discipline
de chaque méthode de service. Le renforcement de cette barrière est traité au
chapitre~\ref{chap:perspectives}.
\end{garantie}

\section{Sécurité : JWT et RBAC}

L'authentification repose sur un couple email / mot de passe chiffré par BCrypt. Le token
JWT, signé en HMAC-SHA, porte les rôles et le \texttt{tenantId}. Le contrôle d'accès est
basé sur les rôles (\texttt{USER}, \texttt{ADMIN}, \texttt{PROVIDER}, \texttt{CONSUMER}),
un même utilisateur pouvant cumuler plusieurs rôles. Spring Security joue ici le rôle d'un
pare-feu applicatif filtrant les requêtes selon l'analogie du tableau~\ref{tab:correspondance}.

\section{Persistance et migrations}

Le schéma est géré par Liquibase, garantissant des migrations versionnées et reproductibles
(création du schéma, ajout des tables \texttt{notifications}, \texttt{messages},
\texttt{audit\_events}, colonnes \texttt{tenant\_id} contraintes par clé étrangère et
indexées). Le paramètre \texttt{ddl-auto=none} laisse Liquibase seul maître du schéma en
exécution.

\section{Temps réel et événementiel}

La messagerie instantanée sur les demandes de service est réalisée par \textbf{Redis
publish/subscribe} : le service de messages publie sur un canal
\texttt{messages:<requestId>}. Un \textbf{bus d'événements Kafka} est configuré
(\texttt{KafkaTemplate} côté producteur, \texttt{@KafkaListener} sur le topic
\texttt{bizcore.events} côté consommateur) ; l'infrastructure est posée et le consommateur
opérationnel, la \emph{production} d'événements restant à câbler (perspective au
chapitre~\ref{chap:perspectives}). Le cache Redis est de même configuré (gestionnaire de
cache, TTL par espace) et prêt à être activé par annotations.

\section{Documentation et conteneurisation}

L'API est documentée automatiquement via SpringDoc OpenAPI et explorable par Swagger UI
(annexe~\ref{ann:swagger}). Un \texttt{docker-compose} décrit l'environnement local complet
(application, PostgreSQL, Redis, Kafka avec \emph{healthchecks} et dépendances
ordonnancées) ; en production, l'API et sa base s'exécutent sur Render.
